\begin{figure}[h!]
    \centering
    \resizebox{0.5\textwidth}{!}{
    \begin{tikzpicture}[
        % GLOBAL STYLES
        font=\small,
        node distance=1.5cm and 3cm, % Vertical and Horizontal spacing
        >=Stealth, % Arrow tip style
        thick,
        % NODE STYLES
        idealbox/.style={
            rectangle, 
            draw=blue!30!white, 
            fill=blue!5!white, 
            rounded corners, 
            text width=5.2cm, 
            align=center, 
            inner sep=10pt, 
            drop shadow
        },
        realbox/.style={
            rectangle, 
            draw=red!30!white, 
            fill=red!5!white, 
            rounded corners, 
            text width=5.2cm, 
            align=center, 
            inner sep=10pt, 
            drop shadow
        },
        corebox/.style={
            rectangle, 
            draw=black!60, 
            fill=yellow!10!white, 
            rounded corners=2pt, 
            double, 
            text width=5.2cm, 
            align=center, 
            inner sep=10pt, 
            drop shadow
        },
        plaintext/.style={
            text width=5.2cm, 
            align=center, 
            font=\itshape, 
            color=black!70
        },
        flowbox/.style={
            rectangle,
            draw=black!30,
            fill=black!2,
            rounded corners,
            text width=5.2cm,
            align=center,
            inner sep=10pt,
            drop shadow
        },
        % ARROW STYLES
        bluearrow/.style={->, double, draw=blue!70!black, double distance=1pt},
        redarrow/.style={->, ultra thick, draw=red!70!black}
    ]

    %% --- NODES ---

    % Top Left: The Platonic Ideal
    \node[idealbox] (IDEAL) {
        \textbf{\textcolor{blue!70!black}{The Platonic Ideal}}\\
        $\realrnn$\\
        \footnotesize (Continuous Arithmetic)
    };

    % Top Right: Continuous Trajectories
    % \node[plaintext, right=of IDEAL] (CONTINUOUS) {
    %     Continuous Trajectories on $\mathbb{R}^N$
    % };
    \node[flowbox, right=of IDEAL] (CONTINUOUS) {%
  Continuous Trajectories on $\R^\hiddim$\\\phantom{A}\\
  \tikz{
    \draw[black!50, line cap=round, smooth, samples=120]
      plot[domain=0:3.6] (\x,{1.2mm*sin(2.4*\x r)});
  }
};

    % Bottom Left: The Digital Reality
    \node[realbox, below=of IDEAL] (REAL) {
        \textbf{\textcolor{red!70!black}{The Digital Reality}}\\
         $\finrnn$\\
        \footnotesize (Discrete States)
    };

    % Bottom Right: The Algebraic Core (aligned with REAL)
    \node[corebox, right=of REAL] (CORE) {
        \textbf{The Algebraic Form}\\
        \footnotesize Behavior goverened by \\
        \large $\rnnsemigroup \leq \layerwreath$
    };

    %% --- ARROWS ---

    % Arrow 1: Ideal -> Continuous
    \draw[->, blue!70!black] (IDEAL) -- node[above, font=\small, align=center] {Theoretical\\Dynamics} (CONTINUOUS);

    % Arrow 2: Ideal -> Real (The Functor)
    \draw[bluearrow] (IDEAL) -- node[right, text width=3cm, align=left, font=\small, color=blue!70!black] {$\functor$ (Discretization)} (REAL);

    % Arrow 3: Continuous -> Core (Collapse)
    \draw[->, dashed, color=black!50]
    (CONTINUOUS.south) -- node[right, text width=3.5cm, align=left, font=\footnotesize]
    {Structural Collapse\\(Calculus fails)}
    (CONTINUOUS.south |- CORE.north);

    % Arrow 4: Real -> Core (Realization)
    \draw[redarrow] (REAL) -- node[above, font=\small, color=red!70!black, align=center] {Algebraic\\Realization} (CORE);
    \end{tikzpicture}
    }
    \caption{A structural perspective: from continuous idealizations (top) to algebraic structures (bottom).\looseness=-1}
    \label{fig:categorical}
\end{figure}